\section{Three-parameter sign band}
With $|c_2|=|c_{-4}|=1$, $\lambda_i\to0$, $|\alpha_i\beta|\ll1$, and positive $A,B,\nu$, the radius-eliminated leading criterion becomes $S_i=(\Pi/\alpha_i-1)/\nu$, $\Pi=\beta/(4A^2)$. For $\alpha_1=1$, positive radial roots and positive slow torque of particle 1 require $\Pi>1$ and $S_1<1$, hence $1<\Pi<1+\nu$. Negative slow torque of particle 2 requires $S_2>1$, or $\Pi>\alpha_2(1+\nu)$. For $0<\alpha_2<1$, intersection yields
\begin{equation}
\max\{1,\alpha_2(1+\nu)\}<\Pi<1+\nu.
\end{equation}
This is a zero-mean-speed torque sign criterion. A stationary velocity sign additionally requires a root of the full slow torque balance and positive effective angular damping $-\partial_\omega\mathcal T_i>0$ near that root. It does not establish branch stability or full RF attraction.
